\section{Counting the roots}\label{sec:counting}

This section proves the counting part of the Main Theorem. All statements are for the torus realization
$\mathbb T(65/64,5/4)$ of $L_A$ and, for $K$, the sharp torus $\mathbb T^\sharp(129/128,9/8)$; by
\cref{lem:Lreal} the counts are the same in the realization $\Yp$ of~\cite{RT2019}. We write
\[
  \Omega_A=(0,3)\times(-\tfrac14,\tfrac14),\qquad \Omega_B=(0,3)\times(\tfrac14,\tfrac34)
\]
for the A and B bands of \cref{sec:setting-sectors}, truncated at $\Re s=3$, and $N(L,\Omega)$ for the
number of roots in $\Omega$ counted with algebraic multiplicity.

\begin{theorem}[Counts]\label{thm:counts}
\begin{enumerate}[label=\textup{(\roman*)},leftmargin=*]
\item $L(s)$ is invertible for $\Re s>2.9261$ and $K(s)$ is invertible for $\Re s>1.765$, for all $\Im s$.
\item $N(L_A,\Omega_A)=3$ and $N(L_A,\Omega_B)=1$.
\item $L_A$ has no roots on $\partial\Omega_A\cup\partial\Omega_B$ other than $s=0$, which is an algebraically
  simple root.
\end{enumerate}
\end{theorem}

\begin{corollary}\label{cor:four}
Modulo $\tfrac i2\Z$, $L$ has exactly four roots in $\{\Re s>0\}$, counted with algebraic multiplicity. On the
imaginary axis its only root is $s\equiv0$, which is algebraically simple.
\end{corollary}

\begin{proof}
By \cref{lem:sectorB}, the roots of $L$ modulo $\tfrac i2\Z$ correspond, with multiplicities, to the roots of
$L_A$ modulo $i\Z$, which are counted in the strip $\widetilde Q=\{-\tfrac14<\Im s\le\tfrac34\}$. By~(i)
there are none with $\Re s\ge3$, and by~(iii) none on the lines $\Im s=\tfrac14$ and $\Im s=-\tfrac14\equiv\tfrac34$
with $0<\Re s<3$. So the roots with $\Re s>0$ are those in $\Omega_A\cup\Omega_B$, and there are $3+1$ of
them. On the axis, the segment $\{\Re s=0,\ -\tfrac14\le\Im s\le\tfrac34\}$ is a full period of $L_A$ and lies
in $\partial\Omega_A\cup\partial\Omega_B$, so~(iii) applies.
\end{proof}

The proof of \cref{thm:counts}~(i) is in \cref{sec:counting-large}. Parts (ii) and (iii) are proved in
\cref{sec:counting-rouche,sec:counting-reference,sec:counting-windows,sec:counting-finite} by an
operator version of Rouch\'e's theorem: after a rank-two deflation of the roots $0$ and $\mu$, $L_A$ is
compared on $\partial\Omega$ with a reference whose zeros are the eigenvalues of a $14016\times14016$
matrix, and these are counted from a Schur decomposition with certified residual.

\subsection{Large real part}\label{sec:counting-large}

On the torus space the background operators are pointwise. With $w=e^{i\tau/2}$ and $z=e^{i\theta}$ on the
torus $|w|=a$, $|z|=b$, multiplication by a field $f$ becomes multiplication by its analytic extension
$f(w,z)$, the reflection $P$ becomes $z\mapsto-z$, and the division by $\xi$ becomes multiplication by
$1/\xi(z)$, $\xi(z)=(z+z^{-1})/2$, which is bounded since $b>1$. Consequently
$\Re\langle -Bx,x\rangle$ is the integral over the torus of a Hermitian form in the seven values
$(x_1(z),x_2(\pm z),x_3(\pm z),x_4(\pm z))$ (using $x_1(-z)=-x_1(z)$), with coefficients given by the background,
and
\[
  \max\Re W(-B)\le\sup_{\text{torus}}\lambda_{\max}\big(N^{-1/2}\operatorname{Herm}(-\mathcal S)N^{-1/2}\big),
\]
where $\mathcal S$ is the $7\times7$ symbol and $N$ accounts for the multiplicity of the values. Here
$W(T)=\{\langle Tx,x\rangle:\|x\|=1\}$ is the numerical range. This bound captures all cancellations
between components and Fourier phases.

\begin{lemma}[Numerical range of the background]\label{lem:F3}
In $\mathbb T(65/64,5/4)$, $\max\Re W(-B)\le2.9492$ and $\|B\|\le3.9642$; for the free part,
$\max\Re W(-J)\le-0.02319$. In $\mathbb T^\sharp(129/128,9/8)$, $\max\Re W(-B^\sharp)\le1.4396$ and
$\|B^\sharp\|\le1.6497$, and $\max\Re W(-K(0))\le 1.765$.
\end{lemma}

\begin{proof}[Method of proof]
The supremum of $\lambda_{\max}$ of the symbol over the torus is bounded on a grid of $2048\times4096$
centres for $L$ and $1024\times2048$ for $K$. At each centre the symbol is evaluated in ball arithmetic (Arb), and the inequality is verified by a
$LDL^{\mathsf T}$ factorization in balls of the shifted matrix. Between centres, Lipschitz constants of the
background fields, computed in exact rational arithmetic, give the slack, also evaluated in ball arithmetic.
The Lipschitz slack added between centres is at most $0.0772$ for $L$ and $0.0536$ for $K$. The free part is handled by a closed
form. Details and file references are in \cref{sec:computation}.
\end{proof}

\begin{proof}[Proof of \cref{thm:counts}~(i)]
For $x$ in the domain, $\Re\langle L(s)x,x\rangle\ge(\Re s-2.9492+0.02319)\|x\|^2$, so
$\|L(s)x\|\ge(\Re s-2.9261)\|x\|$, and $L(s)$ is injective for $\Re s>2.9261$. Since $L(s)$ is Fredholm of
index zero (\cref{prop:fredholm}), it is invertible. The argument for $K$ is the same.
\end{proof}

The bound $2.9261$ should be compared with the largest root, $\sphys\approx0.733$: the exclusion is not sharp,
which is why the rectangles $\Omega_A,\Omega_B$ extend to $\Re s=3$. Before the symbol was used, block-by-block
majorants gave a bound larger by two orders of magnitude.

\subsection{Rouch\'e's theorem for operator families, and deflation}\label{sec:counting-rouche}

We use the operator version of Rouch\'e's theorem of Gohberg and Sigal~\cite{GohbergSigal1971}, in its
homotopy form. Let $\Omega$ be a bounded domain with piecewise smooth boundary, and $H,R$ analytic families of
Fredholm operators of index zero on a neighbourhood of $\overline\Omega$. Suppose moreover that $H-R$ is compact. If
$R+t(H-R)$ is invertible on $\partial\Omega$ for every $t\in[0,1]$, then $H$ and $R$ have the same number of zeros in $\Omega$, counted with
algebraic multiplicity. A sufficient condition is $\|R(s)^{-1}(H(s)-R(s))\|<1$ on $\partial\Omega$.

\emph{Normalization.} The free part $J+s$ is triangular in the Chebyshev index within each Fourier mode, with
diagonal $\mu(n+1)+s+im/2$, so $Q_0(s):=(J+s)^{-1}$ is analytic and invertible in $\Re s>-\mu$. The family
$H(s)=\widetilde L(s)Q_0(s)$ has the same zeros as $\widetilde L$, with the same multiplicities, and
$H(s)=I+(B+E)Q_0(s)$, with $E=\widetilde L-L$ the deflation below, is the identity plus a compact operator.

\emph{Deflation.} The roots $0$ and $\mu$ are known exactly: $L(s)x_0=s\,x_0$ for $x_0=\partial_\tau\omegabar$
(translation invariance in $\tau$), and $L(s)x_1=(s-\mu)\,x_1$ for $x_1=g_*$ (\cref{sec:physical-defs}). The
root $0$ lies on $\partial\Omega_A$, and the root $\mu$ keeps the resolvent large between $0$ and $\mu$. We
remove both by a finite-rank perturbation.

\begin{lemma}[Brauer deflation]\label{lem:brauer}
Let $\phi_0,\phi_1$ be bounded functionals and $\delta_0,\delta_1\in\C$, and set
$\widetilde L(s)=L(s)+\delta_0x_0\phi_0^*+\delta_1x_1\phi_1^*$. Then
\[
  \det\big(I+L(s)^{-1}(\widetilde L(s)-L(s))\big)=\frac{q(s)}{s(s-\mu)},
\]
where $q$ is the quadratic polynomial $\det\big(\operatorname{diag}(s,s-\mu)+[\delta_j\phi_i^*x_j]_{ij}\big)$, and
for every $s_0$ the algebraic multiplicities satisfy
\[
  \mult_{\widetilde L}(s_0)=\mult_L(s_0)+\operatorname{ord}_{s_0}\frac{q(s)}{s(s-\mu)} .
\]
\end{lemma}

\begin{proof}
Since $L(s)^{-1}x_0=x_0/s$ and $L(s)^{-1}x_1=x_1/(s-\mu)$, the perturbation $L(s)^{-1}(\widetilde L-L)$ has
rank two with range in $\operatorname{span}\{x_0,x_1\}$, and the determinant reduces to a $2\times2$
determinant. The multiplicity formula is the logarithmic-residue form of the Gohberg--Sigal theory for the
factorization $\widetilde L=L\,(I+L^{-1}(\widetilde L-L))$, in which the second factor is a finite-rank
meromorphic perturbation of the identity whose logarithmic multiplicity is the order of its determinant.
\end{proof}

We take $\phi_i$ supported in the finite box $Z$ below, biorthogonal to explicit approximations $x_j^A$ computed
from the data $\omega_A$ ($\phi_i^*x_j^A=\delta_{ij}$), and $(\delta_0,\delta_1)=(1,1+\mu_A)$. With
$e_{ij}=\phi_i^*(x_j-x_j^A)$,
\[
  q(s)=(s+1+e_{00})\big(s-\mu+\delta_1(1+e_{11})\big)-\delta_0\delta_1e_{01}e_{10},
\]
so the two deflated roots move to approximately $-1$. The errors satisfy $\|P_Z(x_0-x_0^A)\|\le1.5\cdot10^{-9}$
and $\|x_1-x_1^A\|\le2.51\cdot10^{-8}$, so $|e_{ij}|\le\|\phi_i\|\,\|x_j-x_j^A\|\le5.51\cdot2.51\cdot10^{-8}$, and both factors of $q$ have modulus
$\ge1-10^{-6}$ in $\Re s\ge0$. Hence $q$ has no zeros in $\Re s\ge0$, and by \cref{lem:brauer}
\begin{equation}\label{eq:brauer-count}
  N(L_A,\Omega_A)=N(\widetilde L_A,\Omega_A)+1,\qquad N(L_A,\Omega_B)=N(\widetilde L_A,\Omega_B),
\end{equation}
the $+1$ coming from the pole of $q(s)/(s(s-\mu))$ at $\mu\in\Omega_A$. Moreover, at every $s$ with $\Re s\ge0$ where
$\widetilde L_A(s)$ is invertible, $\mult_L(s)=1$ if $s\in\{0,\mu\}$ and $\mult_L(s)=0$ otherwise. In the certificates the
deflation uses $x_j^A$, and the differences $x_j-x_j^A$ are carried as perturbations.

\subsection{The reference and its verification on the contour}\label{sec:counting-reference}

Partition the coefficient space of sector A into
\begin{itemize}[leftmargin=*]
\item the finite box $Z=\{|m|\le31,\ 0\le n<128\}$, of dimension $14016$;
\item $26$ \emph{windows} $J$ of at most $16$ consecutive Fourier modes, covering the rest of the box
  $\{|m|<200,\ n<400\}$;
\item the \emph{far} part, the complement.
\end{itemize}
Let $P_Z,P_J,P_{\mathrm{far}}$ be the coordinate projections and $\widehat H_{ab}(s)=P_aH(s)P_b$. The reference is the
block-diagonal family
\[
  R(s)=\operatorname{diag}\big(\widehat H_{ZZ}(s),\ \widehat H_{JJ}(s)\ (J\text{ windows}),\ I_{\mathrm{far}}\big).
\]
It is analytic in $\Re s>-\mu$, and its zeros in $\Omega$ are those of the blocks. In the certificates the finite
block is taken one step further from the operator: with $W$ the diagonal matrix of the weights, $\widehat H_{ZZ}$ is
replaced by $R_Z(s)=W^{-1}(\mathsf A+s)W\,(P_Z(J+s)P_Z)^{-1}$, where $\mathsf A$ is the stored floating-point matrix that
approximates $W\widetilde L_Z(0)W^{-1}$, built from the data $\omega_A$ and the vectors $x_j^A$ and taken as an exact
matrix. The difference $P_Z\widetilde LP_Z-W^{-1}\mathsf AW$, due to the background error, the deflation vectors and
the rounding in the assembly, is part of $H-R$ and is bounded in the entries below.

\begin{lemma}[The finite block]\label{lem:ZZ}
$\widehat H_{ZZ}(s)=\widetilde L_Z(s)\,(P_Z(J+s)P_Z)^{-1}$, where $\widetilde L_Z(s)=P_Z\widetilde L(s)P_Z$ is the
$14016\times14016$ truncation. In particular the zeros of $\widehat H_{ZZ}$ in $\Omega$, with multiplicities, are
the points $s$ with $-s$ an eigenvalue of the matrix $\widetilde L_Z(0)$, counted with algebraic multiplicity. The same
holds for $R_Z$, with $\mathsf A$ in place of $W\widetilde L_Z(0)W^{-1}$.
\end{lemma}

\begin{proof}
$J+s$ preserves the Fourier index and is upper triangular in $n$, so the range of $P_Z$ is invariant under
$J+s$ and under $Q_0(s)$: $Q_0(s)P_Z=P_ZQ_0(s)P_Z=(P_Z(J+s)P_Z)^{-1}$. Hence
$\widehat H_{ZZ}(s)=P_Z\widetilde L(s)Q_0(s)P_Z=P_Z\widetilde L(s)P_Z\,(P_Z(J+s)P_Z)^{-1}$, and the second factor is
invertible and analytic in $\Re s>-\mu$.
\end{proof}

\emph{The Perron criterion.} To verify that $R+t(H-R)$ is invertible on $\partial\Omega$ for all $t\in[0,1]$, we
use approximate inverses of the diagonal blocks, $V_Z(s)\approx\widehat H_{ZZ}(s)^{-1}$, $V_J\approx\widehat
H_{JJ}(c)^{-1}$ at a fixed centre $c$, and $I$ for the far part, and bound all block entries.

\begin{lemma}[Perron criterion]\label{lem:perron}
Let $X=\bigoplus_iX_i$, $A=R+E$ with $R=\operatorname{diag}(R_i)$, and let $V_i$ be bounded operators with
$\|I-V_iR_i\|\le\rho_i$ and $\|V_iE_{ij}\|\le n_{ij}$. If the spectral radius $\theta$ of the nonnegative matrix
$\operatorname{diag}(\rho_i)+(n_{ij})$ is $<1$, then $R+tE$ is injective for all $t\in[0,1]$; if $R+tE$ is Fredholm of
index zero, it is invertible.
\end{lemma}

\begin{proof}
With $V=\operatorname{diag}(V_i)$, $V(R+tE)=I-\big((I-VR)-tVE\big)$. In the norm
$\max_iw_i^{-1}\|x_i\|$, with $w$ a positive Perron vector of a slightly larger matrix, the operator in
parentheses has norm $\le\theta+\epsilon<1$, since its block norms are entrywise bounded by the matrix above
(and are increasing in $t$). So $V(R+tE)$ is invertible, and $R+tE$ is injective.
\end{proof}

The levels are $Z$, the windows and the far part; this is the \emph{three-level Perron} bound. The entries are
computed as follows (\cref{sec:computation}).
\begin{itemize}[leftmargin=*]
\item \emph{Level $Z$.} $V_Z(s)$ is applied through a Schur decomposition of $\mathsf A$, used
  only as a solver; residuals are certified a posteriori. A bound $t_p\ge\|(\mathsf A+p)^{-1}\|$ is certified by a Loewner-order inequality $t^2MM^*-I\succeq0$ (with preconditioning near the
  B-band root). The coupling $Z\leftarrow$ tail is factored as a low-rank part $F\Phi$ with a certified
  remainder; with the scaling $D_Z=J_Z(c)Q_{ZZ}(s)$, the entry tail$\,\leftarrow Z$ becomes independent of $s$.
\item \emph{Windows.} Each $V_J$ is the exact inverse of the computed block at the centre, and the blocks are
  bounded by block-Jacobi estimates; the dependence on $s$ is bounded by $\rho_J(s)\le\rho_J(c)+d\,K_J/(1-d\,n_{Q,J})$,
  $d=|s-c|$, with constants computed at the centre (and a refined formula for the central windows).
\item \emph{Far part.} Closed forms for $\|Q_0(s)P_{\mathrm{far}}\|$, uniform on discs.
\item \emph{Background and deflation errors.} The difference between $\omegabar$ and the data, and the errors
  $x_j-x_j^A$, enter as global perturbations of all entries.
\end{itemize}

\emph{Discs.} The bound of \cref{lem:perron} is made uniform on a disc around each point $p$ of the contour:
the variation of the level-$Z$ entries is controlled to first order by structured quantities computed by Schur
solves, and to second order by the resolvent bound, using that $\sigma_{\min}(\widetilde L_Z+s)$ is
$1$-Lipschitz in $s$. For each disc the radius is chosen by bisection on a floating-point estimate of $\theta$, with target $0.999$; the
certified rational bounds are then computed for the chosen radius, and may exceed the target slightly, as long as they
stay below $1$. The contour of $\Omega_A$ is
covered by $96$ discs, with $\max\theta\le0.99842$, and that of $\Omega_B$ by the $96$ discs of $\Omega_A$ on its
lower edge and $78$ further discs, with $\max\theta\le0.99908$; the radii range from $0.0028$ near the B-band
root to $0.4995$ at $\Re s\ge2$. Coverage of the four edges by the union of the discs is checked in rational
arithmetic. This proves that $R+t(H-R)$ is invertible on $\partial\Omega_A$ and on $\partial\Omega_B$, hence
\[
  N(\widetilde L_A,\Omega)=N(R,\Omega)\qquad(\Omega=\Omega_A,\Omega_B),
\]
and that $\widetilde L_A$ is invertible on $\partial\Omega_A\cup\partial\Omega_B$. By the remark after
\eqref{eq:brauer-count}, this is \cref{thm:counts}~(iii).

\subsection{The windows have no zeros}\label{sec:counting-windows}

The Perron criterion on the contour gives $N(\widetilde L,\Omega)=N(R,\Omega)=N(\widehat H_{ZZ},\Omega)+
\sum_JN(\widehat H_{JJ},\Omega)$, but says nothing about zeros of the window blocks inside $\Omega$.

\begin{lemma}\label{lem:windows}
For every window $J$, $\widehat H_{JJ}(s)$ is invertible for all $s\in\overline{\Omega_A}\cup\overline{\Omega_B}$.
\end{lemma}

\begin{proof}
Partition each band into $10$ rectangles $\Omega_k$, one for each tail centre $c_k$. On $\Omega_k$ the function
$F_k(s)=I-V_J(c_k)\widehat H_{JJ}(s)$ is analytic, so $\|F_k(s)\|$ is subharmonic and attains its maximum on
$\partial\Omega_k$. On the boundary, $\|F_k\|$ is bounded by the same sensitivity estimates as in the discs,
evaluated on discs around sample points of the boundary. The maximum is $0.1488$ in the A band and $0.1517$ in the B band.
So $\|F_k\|<1$ on $\Omega_k$, and $\widehat H_{JJ}(s)$ is invertible.
\end{proof}

\subsection{The finite count}\label{sec:counting-finite}

It remains to count the eigenvalues of the stored matrix $\mathsf A$ in $-\Omega$. A floating-point Schur decomposition $\mathsf A\approx UTU^*$ gives the matrix $\mathsf B=UTU^{-1}$,
whose eigenvalues are exactly the diagonal entries of $T$ (with $U$ invertible).

\begin{lemma}[Finite count]\label{lem:schur}
Let $\varepsilon_B\ge\|\mathsf A-\mathsf B\|$ and $t(s)\ge\|(\mathsf A+s)^{-1}\|$ on $\partial\Omega$. If
$\sup_{\partial\Omega}t(s)\,\varepsilon_B<1$, then $\mathsf A$ has exactly $\#\{i:-T_{ii}\in\Omega\}$ eigenvalues
in $-\Omega$, with multiplicity.
\end{lemma}

\begin{proof}
$\mathsf A+s+t'(\mathsf B-\mathsf A)$ is invertible on $\partial\Omega$ for $t'\in[0,1]$, by a Neumann series, so the
number of eigenvalues inside does not change along the homotopy.
\end{proof}

We have $\varepsilon_B\le\|\mathsf AU-UT\|_F/\sqrt{1-\|I-U^*U\|_F}$ with $\|\mathsf AU-UT\|_F\le5.97\cdot10^{-5}$ and
$\|I-U^*U\|_F\le1.31\cdot10^{-6}$, and $t(s)\le t_p/(1-rt_p)$ on each disc, so $\sup t\le559.2$ and
$t\,\varepsilon_B\le0.0334$. The diagonal of $T$ has exactly two entries $-T_{ii}$ in $\Omega_A$, at
$0.401025\ldots$ and $0.733180\ldots$, and exactly one in $\Omega_B$, at $0.0414194\ldots+0.5i$; the nearest
entry outside $\Omega_A$ is at distance $0.168$ from its boundary. Hence $N(\widehat H_{ZZ},\Omega_A)=2$ and
$N(\widehat H_{ZZ},\Omega_B)=1$.

\begin{proof}[Proof of \cref{thm:counts}~(ii)]
By \cref{lem:windows,lem:schur,lem:ZZ}, $N(R,\Omega_A)=2$ and $N(R,\Omega_B)=1$. By the Perron verification on the
contours, $N(\widetilde L_A,\Omega)=N(R,\Omega)$, and by~\eqref{eq:brauer-count},
$N(L_A,\Omega_A)=2+1=3$ and $N(L_A,\Omega_B)=1$.
\end{proof}

The positions $0.401$, $0.733$ and $0.0414+0.5i$ of the finite reference are not used in the count; the roots of
$L$ are localized independently in \cref{sec:roots}. The Rouch\'e argument only guarantees that there are three
roots in $\Omega_A$ and one in $\Omega_B$.
